\documentclass[a4paper,11pt]{article}
\usepackage[a4paper, top=2.5cm, bottom=2.5cm, left=2.5cm, right=2.5cm]{geometry}
\usepackage[utf8]{inputenc}
\usepackage[T1]{fontenc}
\usepackage[ngerman]{babel}
\usepackage{microtype}
\usepackage{amsmath,amssymb}
\usepackage{booktabs}
\usepackage{array}
\usepackage{hyperref}
\usepackage{newtxtext,newtxmath}
\usepackage{graphicx}
\usepackage{parskip}
\usepackage{fancyhdr}
\usepackage{titlesec}

\hypersetup{
  pdftitle={Kooperationsvorschlag RFT -- Prof. Jing Zhang},
  pdfauthor={Dominic Ren\'e Schu},
  colorlinks=true,
  linkcolor=black,
  urlcolor=blue,
  citecolor=black
}

\pagestyle{fancy}
\fancyhf{}
\lhead{\small Resonanzfeldtheorie -- Kooperationsvorschlag}
\rhead{\small Dominic Ren\'e Schu}
\cfoot{\thepage}
\renewcommand{\headrulewidth}{0.4pt}

\titleformat{\section}{\large\bfseries}{}{0em}{}[\titlerule]
\titleformat{\subsection}{\normalsize\bfseries}{}{0em}{}

\setlength{\parindent}{0pt}

%% ====================================================================
\begin{document}

%% ---- Titelseite ----
\thispagestyle{empty}
\begin{center}
  {\LARGE\bfseries Kooperationsvorschlag}\\[0.5em]
  {\large\bfseries Falsifizierbare Vorhersage nichtlinearer\\
  Quantenmechanik in $^{87}$Rb-Bose-Einstein-Kondensaten}\\[2em]
  {\large Dominic Ren\'e Schu}\\[0.3em]
  {\normalsize Resonanzfeldtheorie (RFT)\\
  \href{https://github.com/DominicReneSchu/RFT}{github.com/DominicReneSchu/RFT}}\\[2em]
  \rule{0.6\textwidth}{0.4pt}\\[1em]
  {\normalsize Adressat:}\\[0.3em]
  {\large\bfseries Professor Jing Zhang (张靖)}\\
  {\normalsize State Key Laboratory of Quantum Optics and Quantum Optics Devices\\
  Shanxi University, Taiyuan, China}\\[2em]
  {\normalsize \today}
\end{center}

\vfill
\begin{center}
\small
Dieses Dokument ist ein formeller wissenschaftlicher Kooperationsvorschlag.
Die vollständige theoretische Grundlage, alle Simulationen und der
experimentelle Vorschlag sind öffentlich zugänglich unter:\\
\href{https://github.com/DominicReneSchu/RFT}{https://github.com/DominicReneSchu/RFT}
\end{center}
\newpage

%% ---- Zusammenfassung ----
\section{Zusammenfassung}

Die \textbf{Resonanzfeldtheorie} (RFT) ist ein axiomatisches Framework, das
die Schrödinger-Gleichung als Spezialfall einer allgemeineren, nichtlinearen
Dynamik herleitet. Über ein Dichte-Rückkopplungsmodell
\[
  \dot{\Delta\varphi} = \lambda \int |\psi|^4\,\mathrm{d}x
\]
erzeugt die RFT im Grenzfall $\lambda \to 0$ exakt die Standard-Quantenmechanik,
macht für $\lambda > 0$ jedoch messbar abweichende Vorhersagen.

Der vorliegende Kooperationsvorschlag legt einen vollständig ausgearbeiteten,
\textbf{falsifizierbaren experimentellen Vorschlag} vor:
\[
  \boxed{|\Delta\langle x\rangle|_{\mathrm{SI}} \approx 2.0 \cdot \lambda\;\mu\mathrm{m}}
\]
Die Vorhersage basiert auf numerisch verifizierten Skalierungsgesetzen,
vollständiger SI-Kalibrierung für $^{87}$Rb in einer harmonischen Falle
und einem sauber formulierten statistischen Testprotokoll.

\textbf{Ziel der Kooperation:} Experimentelle Überprüfung der Vorhersage mit
der vorhandenen BEC-Infrastruktur der Gruppe Zhang (Shanxi University).
Das Experiment ist ein \emph{Bound-Experiment}: Es liefert entweder einen
Nachweis des RFT-Effekts oder setzt eine obere Schranke für $\lambda$ ---
beides ist ein wissenschaftlich wertvolles Ergebnis.

\newpage
\tableofcontents
\newpage

%% ====================================================================
\section{Theoretischer Hintergrund}

\subsection{Resonanzfeldtheorie: Axiom 4 und Schrödinger-Gleichung}

Die RFT ist auf sieben Axiome aufgebaut. Für das vorliegende Experiment
ist \textbf{Axiom 4} (Kopplungseffizienz) zentral: Die effektive Kopplung
zwischen einem quantenmechanischen System und dem Resonanzfeld wird durch
die Phasendifferenz $\Delta\varphi$ moduliert:
\[
  \varepsilon(\Delta\varphi) = \cos^2\!\left(\frac{\Delta\varphi}{2}\right) \in [0, 1].
\]
Der Resonanz-Hamiltonoperator lautet:
\[
  \hat{H}_{\mathrm{res}} = \hat{H}_0 + \varepsilon(\Delta\varphi)\,\hat{V}_{\mathrm{Kopplung}}.
\]
Für konstantes $\Delta\varphi$ (statischer Grenzfall) folgt daraus exakt die
Standard-Schrödinger-Gleichung --- das Korrespondenzprinzip ist damit analytisch
und numerisch belegt (Fidelity $= 1{,}000000000000$ für alle getesteten $\Delta\varphi$).

\subsection{Dichte-Rückkopplungsmodell}

Für die dynamische Erweiterung wird $\Delta\varphi(t)$ zu einem zeitabhängigen
Feld, das an die Wellenfunktion rückkoppelt:
\[
  \Delta\varphi(t + \mathrm{d}t) = \Delta\varphi(t)
  + \lambda \cdot \underbrace{\int |\psi(x,t)|^4\,\mathrm{d}x}_{F[\psi]} \cdot \mathrm{d}t.
\]
Dieses Modell ist aus dem RFT-Wirkungsfunktional
\[
  S[\psi, \Delta\varphi] = \int \mathrm{d}t \left[
    \langle\psi|i\hbar\partial_t - \hat{H}_0|\psi\rangle
    - \varepsilon(\Delta\varphi)\langle V\rangle_\psi
    + \frac{\mu}{2}(\dot{\Delta\varphi})^2
  \right]
\]
über Euler-Lagrange-Gleichungen herleitbar (überdämpfter Grenzfall).
Die Rückkopplung $\psi \to \Delta\varphi \to V_{\mathrm{eff}} \to \psi$ ist
strikt lokal und verletzt das No-Signaling-Prinzip nicht.

\subsection{Störungstheorie und Skalierungsgesetze}

Für $\lambda \ll 1$ liefert die Störungstheorie 1.~Ordnung analytische
Ausdrücke für alle Observablen. Die Numerik bestätigt diese bis auf einen
relativen Fehler von $1{,}3 \times 10^{-4}$. Die zentralen Skalierungsgesetze:

\begin{table}[h]
\centering
\caption{Skalierungsexponenten der RFT-Observablen (Störungstheorie vs.\ Numerik)}
\begin{tabular}{lcccr}
\toprule
Observable & Exp.\ (numerisch) & Exp.\ (Theorie) & Abweichung \\
\midrule
$1 - F$ (Fidelity) & $2{,}001$ & $2$ & $0{,}05\,\%$ \\
$|\Delta\langle x\rangle|$ & $1{,}001$ & $1$ & $0{,}1\,\%$ \\
$|\Delta\langle p\rangle|$ & $1{,}001$ & $1$ & $0{,}1\,\%$ \\
$\max|\Delta\psi|$ & $0{,}999$ & $1$ & $0{,}1\,\%$ \\
\bottomrule
\end{tabular}
\end{table}

Die numerisch bestimmten Präfaktoren sind:
\[
  |\Delta\langle x\rangle| \approx 4{,}9 \cdot \lambda, \quad
  |\Delta\langle p\rangle| \approx 3{,}2 \cdot \lambda, \quad
  1 - F \approx 138 \cdot \lambda^2.
\]
Dass $1 - F \sim \lambda^2$ gilt, folgt direkt aus
$|\langle\psi_0|\psi_0 + \lambda\psi_1\rangle|^2
= 1 - \lambda^2\|\psi_1\|^2 + \mathcal{O}(\lambda^3)$.

%% ====================================================================
\section{SI-Kalibrierung}

\subsection{Physikalisches System}

Das Referenzsystem sind \textbf{ultrakalte $^{87}$Rb-Atome in einer
harmonischen Falle} --- das weltweit etablierteste System für
Präzisions-Quantenmechanik:

\begin{table}[h]
\centering
\caption{Systemparameter für $^{87}$Rb}
\begin{tabular}{llr}
\toprule
Parameter & Symbol & Wert \\
\midrule
Atommasse & $m$ & $86{,}909\,\mathrm{u} = 1{,}443 \times 10^{-25}\,\mathrm{kg}$ \\
Fallenfrequenz & $\omega$ & $2\pi \times 100\,\mathrm{Hz}$ (einstellbar) \\
Planck-Konstante & $\hbar$ & $1{,}055 \times 10^{-34}\,\mathrm{J\cdot s}$ \\
\bottomrule
\end{tabular}
\end{table}

\subsection{Vollständige Skalenabbildung}

Die dimensionslose Simulation ($\hbar = m = 1$) wird über drei
Konversionsgrößen auf SI abgebildet. Die Kalibrierung ist durch
6 unabhängige Konsistenzrelationen (z.\,B.\ $\hbar = m\,\ell^2/\tau$,
$E \cdot \tau = \hbar$) überbestimmt und vollständig konsistent.

\begin{table}[h]
\centering
\caption{SI-Kalibrierung für $\omega = 2\pi \times 100\,\mathrm{Hz}$}
\begin{tabular}{llll}
\toprule
Größe & Formel & Wert ($\omega = 2\pi \times 100\,\mathrm{Hz}$) \\
\midrule
Oszillatorlänge & $a_{\mathrm{ho}} = \sqrt{\hbar/(m\omega)}$ & $1{,}08\,\mu\mathrm{m}$ \\
Längeneinheit & $\ell = V_s^{1/4} \cdot a_{\mathrm{ho}}$ & $0{,}41\,\mu\mathrm{m}$ \\
Zeiteinheit & $\tau = \sqrt{V_s}/\omega$ & $0{,}225\,\mathrm{ms}$ \\
Energieeinheit & $E = \hbar/\tau$ & $2{,}92 \times 10^{-12}\,\mathrm{eV}$ \\
Impulseinheit & $p = \hbar/\ell$ & $2{,}60 \times 10^{-28}\,\mathrm{kg\cdot m/s}$ \\
Simulationszeit & $T = 20\,\tau$ & $4{,}5\,\mathrm{ms}$ \\
\bottomrule
\end{tabular}
\end{table}

\subsection{Herleitung der Skalenabbildung}

Die dimensionslose Schrödinger-Gleichung mit $\hbar = m = 1$,
\[
  i\,\partial_{\tilde{t}}\,\psi = -\tfrac{1}{2}\,\partial_{\tilde{x}}^2\,\psi
  + V(\tilde{x})\,\psi,
\]
wird über $x = \ell\,\tilde{x}$, $t = \tau\,\tilde{t}$ auf SI abgebildet.
Das Kopplungspotential $V(\tilde{x}) = \tfrac{1}{2}V_s\,\tilde{x}^2$
wird mit dem physikalischen Potential $V_{\mathrm{phys}}(x) = \tfrac{1}{2}m\omega^2 x^2$
identifiziert:
\[
  V_s = \frac{m^2\omega^2\ell^4}{\hbar^2}
  \quad\Longrightarrow\quad
  \ell = V_s^{1/4} \cdot \sqrt{\frac{\hbar}{m\omega}} = V_s^{1/4} \cdot a_{\mathrm{ho}}.
\]

%% ====================================================================
\section{Falsifizierbare Vorhersage}

\subsection{Zentrale Vorhersage}

Die RFT sagt eine messbare Schwerpunktsverschiebung des Wellenpakts vorher:
\[
  \boxed{|\Delta\langle x\rangle|_{\mathrm{SI}}
  = C_x \cdot \lambda \cdot \ell
  \approx 2{,}0 \cdot \lambda\;\mu\mathrm{m}}
\]
mit $C_x \approx 4{,}9$ (numerisch bestimmt) und
$\ell \approx 0{,}41\,\mu\mathrm{m}$ bei $\omega = 2\pi \times 100\,\mathrm{Hz}$.

\subsection{Vorhersagetabelle}

\begin{table}[h]
\centering
\caption{Vorhersage der Positionsverschiebung und Fidelity-Verlust für verschiedene $\lambda$}
\begin{tabular}{rrrrl}
\toprule
$\lambda$ & $|\Delta\langle x\rangle|$ [$\mu$m] & $1 - F$ & Detektierbar? \\
\midrule
$0{,}001$ & $0{,}002$ & $1{,}4 \times 10^{-4}$ & Nein (Einzelschuss) \\
$0{,}01$  & $0{,}020$ & $0{,}014$               & Ja (10\,000 Schuss) \\
$0{,}05$  & $0{,}099$ & $0{,}345$               & Ja (100 Schuss) \\
$0{,}1$   & $0{,}199$ & $1{,}38$                & Ja (100 Schuss) \\
$0{,}5$   & $0{,}994$ & $34{,}5$                & Ja (Einzelschuss) \\
$1{,}0$   & $1{,}987$ & $138$                   & Ja (Einzelschuss) \\
\bottomrule
\end{tabular}
\end{table}

\subsection{Frequenzabhängigkeit und optimaler Bereich}

Niedrigere Fallenfrequenz ergibt größeres $a_{\mathrm{ho}}$ und damit
bessere Sensitivität:

\begin{table}[h]
\centering
\caption{Sensitivität als Funktion der Fallenfrequenz}
\begin{tabular}{rrrr}
\toprule
$\omega / 2\pi$ [Hz] & $a_{\mathrm{ho}}$ [$\mu$m] & $\ell$ [$\mu$m]
  & $\lambda_{\min}$ (100 Schuss) \\
\midrule
$10$   & $3{,}41$ & $1{,}28$ & $0{,}016$ \\
$50$   & $1{,}53$ & $0{,}57$ & $0{,}036$ \\
$100$  & $1{,}08$ & $0{,}41$ & $0{,}050$ \\
$500$  & $0{,}48$ & $0{,}18$ & $0{,}113$ \\
$1000$ & $0{,}34$ & $0{,}13$ & $0{,}159$ \\
\bottomrule
\end{tabular}
\end{table}

\textbf{Optimaler Bereich:} $\omega \lesssim 2\pi \times 50\,\mathrm{Hz}$
(bei längerer Präparationszeit).

%% ====================================================================
\section{Experimentelles Protokoll}

\subsection{Vorbereitung}

\begin{enumerate}
  \item \textbf{BEC-Präparation:} $^{87}$Rb-Kondensat in magnetischer oder
    optischer Falle, $T < 200\,\mathrm{nK}$, $N \sim 10^5$ Atome.
  \item \textbf{Harmonische Falle:} Einstellbar auf
    $\omega = 2\pi \times (10\text{--}500)\,\mathrm{Hz}$.
    Die axiale Frequenz bestimmt $a_{\mathrm{ho}}$ und damit die Sensitivität.
  \item \textbf{Wellenpaket-Initialisierung:} Impuls-Kick über Bragg-Puls
    oder Raman-Übergang: $\hbar k_0 = \hbar/\ell$.
\end{enumerate}

\subsection{Messung}

\begin{enumerate}
  \setcounter{enumi}{3}
  \item \textbf{Propagation:} Wellenpaket für $t \approx 4{,}5\,\mathrm{ms}$
    propagieren lassen (bei $\omega = 2\pi \times 100\,\mathrm{Hz}$).
  \item \textbf{Absorptionsbildgebung:} Falle abschalten, Time-of-Flight,
    Absorptionsbild aufnehmen. Räumliche Auflösung: ${\sim}\,1\,\mu\mathrm{m}$.
  \item \textbf{Wiederholung:} $N = 100\text{--}10\,000$ identische Durchläufe.
\end{enumerate}

\subsection{Auswertung}

\begin{enumerate}
  \setcounter{enumi}{6}
  \item \textbf{Statistik:} Mittlere Position $\langle x\rangle_{\mathrm{exp}}$
    aus $N$ Messungen.
    Standardfehler: $\sigma_{\Delta x} = \sigma_{\mathrm{single}}/\sqrt{N}$.
  \item \textbf{Hypothesentest:}
    \begin{itemize}
      \item \textbf{Nullhypothese} $H_0$: $\Delta\langle x\rangle = 0$
        (Standard-QM, $\lambda = 0$)
      \item \textbf{Alternativhypothese} $H_1$:
        $|\Delta\langle x\rangle| = C_x \cdot \lambda \cdot \ell > 0$ (RFT)
    \end{itemize}
  \item \textbf{Ergebnis:}
    \begin{itemize}
      \item $|\Delta\langle x\rangle| > 2\sigma$: RFT-Effekt nachgewiesen
        $\to$ $\lambda$ bestimmt
      \item $|\Delta\langle x\rangle| \leq 2\sigma$: Obere Schranke
        $\lambda_{\max}$ gesetzt
    \end{itemize}
\end{enumerate}

%% ====================================================================
\section{Warum die Gruppe von Professor Zhang}

Die Gruppe von Professor Jing Zhang am State Key Laboratory of Quantum
Optics and Quantum Optics Devices (Shanxi University) besitzt präzise die
Infrastruktur und Expertise, die das vorgeschlagene Protokoll erfordert:

\begin{itemize}
  \item \textbf{$^{87}$Rb-BEC-Expertise:} Jahrelange Erfahrung in der
    Präparation und Manipulation von Rubidium-Bose-Einstein-Kondensaten ---
    das exakte System des Protokolls.
  \item \textbf{Absorptionsbildgebung:} Hochauflösende Absorptionsbildgebung
    mit räumlicher Auflösung $\sim 1\,\mu\mathrm{m}$ ist die zentrale
    Messmethode des Vorschlags.
  \item \textbf{Kohärente Dynamik und Wellenpaketmanipulation:}
    Veröffentlichte Arbeiten zur kohärenten Wellenpaketdynamik passen
    direkt zu den Mess\-anforderungen (Impuls-Kick, freie Propagation,
    Positionsmessung).
  \item \textbf{State Key Laboratory Infrastruktur:} Magnetfeldkontrolle,
    einstellbare Fallenfrequenzen und optische Dipolfallen sind im Protokoll
    explizit vorgesehen.
  \item \textbf{Direkte Passung zum Protokoll:} Alle vier Schritte
    (BEC-Präparation, harmonische Falle, Impuls-Kick, Absorptionsbildgebung)
    sind in der Gruppe Zhang Standardprozeduren.
\end{itemize}

%% ====================================================================
\section{Kritische Einordnung}

\subsection{Kohn-Theorem: Unterscheidung von Gross-Pitaevskii-Effekten}

Die RFT-Rückkopplung $\dot{\Delta\varphi} \propto \int |\psi|^4\,\mathrm{d}x$
hat dieselbe funktionale Form wie der Kontaktwechselwirkungsterm der
Gross-Pitaevskii-Gleichung ($g|\psi|^2\psi$). Die konzeptuelle Unterscheidung
liefert das \textbf{Kohn-Theorem}:

In einer \emph{rein harmonischen} Falle ist die Schwerpunktsbewegung
$\langle x\rangle(t)$ exakt harmonisch bei $\omega$ ---
\emph{unabhängig von Atom-Atom-Wechselwirkungen}:
\[
  \text{GP: } \quad \Delta\langle x\rangle = 0 \quad\text{(Kohn-geschützt)}.
\]
Die RFT-Rückkopplung hingegen moduliert die Fallenstärke zeitabhängig
über $\varepsilon(\Delta\varphi(t)) \cdot V$. Diese zeitabhängige Modulierung
\emph{bricht} die Kohn-Bedingung:
\[
  \text{RFT: } \quad |\Delta\langle x\rangle| = C_x \cdot \lambda \cdot \ell \neq 0.
\]
Der GP-Term ändert die \emph{Breite} des Kondensats, nicht den
\emph{Schwerpunkt}. Das RFT-Signal ist daher konzeptuell sauber
von Mean-Field-Effekten getrennt.

\subsection{Systematische Fehler}

\begin{table}[h]
\centering
\caption{Fehlerbudget für $^{87}$Rb-BEC in harmonischer Falle}
\begin{tabular}{lrl}
\toprule
Fehlerquelle & Geschätzte Verschiebung & Skalierung \\
\midrule
GP-Mean-Field auf $\langle x\rangle$ & \textbf{0} (Kohn-Theorem) & nicht anwendbar \\
Potentialanharmonizitäten & ${\sim}\,1\,\mathrm{nm}$ & $\delta\omega/\omega$ \\
Magnetfeldgradienten ($0{,}1\,\mathrm{mG/cm}$) & ${\sim}\,650\,\mathrm{nm}$ & $\mu_B \cdot \nabla B$ \\
Drei-Körper-Verluste & ${\sim}\,11\,\mathrm{nm}$ & $\Delta N / N$ \\
\midrule
\textbf{Gesamt (quadratisch)} & $\mathbf{{\sim}\,650\,\mathrm{nm}}$ & \\
\bottomrule
\end{tabular}
\end{table}

Für $\lambda \lesssim 0{,}3$ dominieren die Magnetfeldgradienten.
Eine Magnetfeldkompensation $< 0{,}01\,\mathrm{mG/cm}$ würde die
Systematik um eine Größenordnung reduzieren.

\subsection{Unterscheidungsprotokolle}

\begin{table}[h]
\centering
\caption{Experimentelle Unterscheidung zwischen GP und RFT}
\begin{tabular}{lll}
\toprule
Protokoll & GP-Abhängigkeit & RFT-Abhängigkeit \\
\midrule
$N$-Scan & Breite $\propto N^{2/5}$ & $\langle x\rangle$-Shift $\neq f(N)$ \\
$a_s$-Scan (Feshbach) & $\propto a_s$ & unabhängig von $a_s$ \\
$\omega$-Scan & Breite $\propto \omega^{-3/5}$ & Shift $\propto \omega^{-1/2}$ \\
$\Delta\varphi_0$-Scan & unabhängig & $\propto \varepsilon'(\Delta\varphi_0)$ \\
\bottomrule
\end{tabular}
\end{table}

\subsection{Das Experiment als Bound-Experiment}

Das vorgeschlagene Experiment ist transparent als \emph{Bound-Experiment}
einzuordnen:
\begin{itemize}
  \item \textbf{Positives Ergebnis:} $|\Delta\langle x\rangle| > 2\sigma$
    $\Rightarrow$ RFT-Effekt nachgewiesen, $\lambda$ quantitativ bestimmt.
  \item \textbf{Null-Ergebnis:} $|\Delta\langle x\rangle| \leq 2\sigma$
    $\Rightarrow$ Obere Schranke $\lambda_{\max}$ gesetzt, RFT-Parameterraum
    signifikant eingeschränkt.
\end{itemize}
Beide Ausgänge sind wissenschaftlich wertvolle, publizierbare Ergebnisse.

%% ====================================================================
\section{Vorgeschlagene Kooperationsstruktur}

\begin{itemize}
  \item \textbf{Theorie und Simulation:}
    Dominic Ren\'e Schu ---
    vollständige analytische Grundlage, numerische Simulationen,
    SI-Kalibrierung, Vorhersageberechnungen.
  \item \textbf{Experiment:}
    Gruppe Zhang (Shanxi University) ---
    $^{87}$Rb-BEC-Präparation, harmonische Falle, Absorptionsbildgebung,
    statistische Auswertung.
  \item \textbf{Gemeinsame Publikation:}
    Beide Parteien als Koautoren in einer peer-reviewed Fachzeitschrift
    (z.\,B.\ \textit{Physical Review Letters} oder
    \textit{Nature Physics}).
  \item \textbf{Offene Wissenschaft:}
    Alle Simulationen, Daten und Auswertungen werden im
    GitHub-Repository öffentlich dokumentiert.
\end{itemize}

%% ====================================================================
\section{Referenzen}

\begin{enumerate}
  \item Dominic Ren\'e Schu, \textit{Resonanzfeldtheorie --- Axiomatisches
    Framework und experimenteller Vorschlag}, GitHub-Repository,
    \url{https://github.com/DominicReneSchu/RFT} (2025/2026).
  \item Vollständiger experimenteller Vorschlag (SI-Kalibrierung, Störungstheorie,
    Falsifizierbare Vorhersage):\\
    \texttt{de/fakten/simulationen/schrödinger/docs/experimental\_proposal.md}
  \item Simulationspaket (Python, numpy):\\
    \texttt{de/fakten/simulationen/schrödinger/python/schrodinger\_1d\_rft\_experiment.py}
  \item Übersicht der Schrödinger-Simulation:\\
    \texttt{de/fakten/simulationen/schrödinger/README.md}
  \item W.~Kohn, \textit{Cyclotron Resonance and de Haas-van Alphen
    Oscillations of an Interacting Electron Gas}, Phys.\ Rev.\ \textbf{123},
    1242 (1961). [Kohn-Theorem]
  \item C.~Pethick und H.~Smith, \textit{Bose-Einstein Condensation in
    Dilute Gases}, Cambridge University Press (2008).
\end{enumerate}

\vspace{2em}
\hrule
\vspace{1em}
\begin{center}
\small
\textit{Dominic Ren\'e Schu --- Resonanzfeldtheorie}\\
\href{https://github.com/DominicReneSchu/RFT}{https://github.com/DominicReneSchu/RFT}
\end{center}

\end{document}
